\begin{explanationbox}
\textbf{先看一个具体交弦。} 对 $x^2/4-y^2=1$ 与 $y=x+2$，交点为
\[
 P\left(-\frac{10}{3},-\frac43\right),\qquad Q(-2,0).
\]
\begin{center}\begin{tikzpicture}[font=\small]
\draw[black!30,rounded corners=2pt] (0,0) rectangle (14.4,8.2);
\node[anchor=west,font=\small\bfseries] at (.5,7.72) {两个根对应两个交点：它们可以在同一支上};
\begin{scope}[shift={(7.2,3.65)},scale=0.95]
\begin{scope}
\clip (-4.5,-2.1) rectangle (4.5,2.1);
\draw[blue,very thick,domain=-2.1:2.1,samples=161,variable=\t] plot ({2*sqrt(1+(\t/1)^2)},\t);
\draw[blue,very thick,domain=-2.1:2.1,samples=161,variable=\t] plot ({-2*sqrt(1+(\t/1)^2)},\t);
\end{scope}
\draw[-Stealth,black!45] (-4.5,0)--(4.5,0) node[right] {$x$};
\draw[-Stealth,black!45] (0,-2.1)--(0,2.1) node[above] {$y$};
\draw[orange!85!black,very thick] (-4.1,-2.1)--(.1,2.1);
\coordinate (P) at (-3.3333333333,-1.3333333333);\fill[black] (P) circle (1.8pt) node[above left=10pt] {$P$};
\coordinate (Q) at (-2.0000000000,0.0000000000);\fill[black] (Q) circle (1.8pt) node[below right=11pt] {$Q$};
\draw[teal!70!black,dashed,thick] (P)--(-3.3333333333,0);\node[above] at (-3.3333333333,0) {$x_1$};\node[above=16pt] at (Q) {$x_2$};\end{scope}
\node[anchor=west] at (.7,.55) {$x^2/4-y^2=1,\quad y=x+2$；这里 $x_1x_2>0$，两点都在左支};
\end{tikzpicture}\end{center}

图中的横坐标投影就是消元方程的两个根。它们都为负，两个交点都在左支；因此\textbf{“双曲线有两支”不等于“直线的两交点分居两支”}。
本例代入只得到 $3x^2+16x+20=0$，不必逐一解根也能读出和与积。
\end{explanationbox}
\clearpage
\begin{explanationbox}
\textbf{二次项非零，只保证可以用韦达。}
固定 $x^2/2-y^2=1$，只改变直线 $y=x+m$ 的截距。
这里 $M=1-2=-1\ne0$，而 $H=m^2-1$ 会变号。
\begin{center}\begin{tikzpicture}[font=\small]
\draw[black!30,rounded corners=2pt] (0,0) rectangle (14.4,8.2);
\node[anchor=west,font=\small\bfseries] at (.5,7.72) {固定方向，只改变截距：无交点、相切、两交点};
\begin{scope}[shift={(7.2,3.65)},scale=0.53]
\begin{scope}
\clip (-7,-4.7) rectangle (7,4.7);
\draw[blue,very thick,domain=-4.7:4.7,samples=161,variable=\t] plot ({1.4142135623730951*sqrt(1+(\t/1)^2)},\t);
\draw[blue,very thick,domain=-4.7:4.7,samples=161,variable=\t] plot ({-1.4142135623730951*sqrt(1+(\t/1)^2)},\t);
\end{scope}
\draw[-Stealth,black!45] (-7,0)--(7,0) node[right] {$x$};
\draw[-Stealth,black!45] (0,-4.7)--(0,4.7) node[above] {$y$};
\draw[red!75!black,densely dashed,thick] (-4.7,-4.7)--(4.7,4.7);
\draw[violet,thick] (-5.7,-4.7)--(3.7,4.7);
\draw[orange!85!black,very thick] (-6.7,-4.7)--(2.7,4.7);
\coordinate (T) at (-2.0000000000,-1.0000000000);\fill[black] (T) circle (1.8pt) ;
\coordinate (P) at (-6.4494897428,-4.4494897428);\fill[black] (P) circle (1.8pt) ;
\coordinate (Q) at (-1.5505102572,0.4494897428);\fill[black] (Q) circle (1.8pt) ;
\end{scope}\node[anchor=west,red!75!black] at (11.1,6.4) {$m=0$：无交点};\node[anchor=west,violet] at (11.1,5.4) {$m=1$：相切};\node[anchor=west,orange!85!black] at (11.1,4.4) {$m=2$：两交点};
\node[anchor=west] at (.7,.55) {$x^2/2-y^2=1,\ y=x+m$；$M=-1$，故比较 $H=m^2-1$ 的符号};
\end{tikzpicture}\end{center}

$m=0$ 时无交点；$m=1$ 时在 $(-2,-1)$ 相切；$m=2$ 时出现两个交点。
这说明\textbf{先检查二次项，再检查实交点}，两步不能合并。
\end{explanationbox}
\clearpage
\begin{explanationbox}
\textbf{二次项为零，发生了什么？}
对 $x^2-y^2=1$，渐近线方向为 $y=\pm x$。沿其中一个方向平移直线：
\begin{center}\begin{tikzpicture}[font=\small]
\draw[black!30,rounded corners=2pt] (0,0) rectangle (14.4,8.2);
\node[anchor=west,font=\small\bfseries] at (.5,7.72) {二次项消失：平行于渐近线，不重合时恰有一个交点};
\begin{scope}[shift={(7.2,3.65)},scale=1.05]
\begin{scope}
\clip (-3.1,-2.1) rectangle (3.1,2.1);
\draw[blue,very thick,domain=-2.1:2.1,samples=161,variable=\t] plot ({1*sqrt(1+(\t/1)^2)},\t);
\draw[blue,very thick,domain=-2.1:2.1,samples=161,variable=\t] plot ({-1*sqrt(1+(\t/1)^2)},\t);
\end{scope}
\draw[-Stealth,black!45] (-3.1,0)--(3.1,0) node[right] {$x$};
\draw[-Stealth,black!45] (0,-2.1)--(0,2.1) node[above] {$y$};
\draw[black!55,densely dashed,thick] (-2.1,-2.1)--(2.1,2.1);
\draw[orange!85!black,very thick] (-3.1,-2.1)--(1.1,2.1);
\coordinate (P) at (-1.0000000000,0.0000000000);\fill[black] (P) circle (1.8pt) node[above left=10pt] {$P$};
\node[anchor=west,black!60] at (1.7,1.2) {$y=x$：渐近线};\end{scope}\node[anchor=west,orange!85!black] at (.8,6.6) {$y=x+1$：一个交点};
\node[anchor=west] at (.7,.55) {$x^2-y^2=(x-y)(x+y)=1$；$y=x+1$ 固定了 $x-y=-1$};
\end{tikzpicture}\end{center}

$y=x+1$ 代入后，$x^2$ 项抵消，只剩 $-2x-1=1$，所以交点恰为 $(-1,0)$。
若去掉截距，$y=x$ 代入得到 $0=1$，因此无交点。
\textbf{这里少一个交点是方程降次，不是出现二重根。}
渐近线是曲线无限延伸时逐渐靠近的直线；图中曲线与它不相交。
\end{explanationbox}
\clearpage
% ===== 新增：两点相减与交叉乘减的直观理解（原内容保持不变）=====
\begin{explanationbox}
\textbf{新增速记：两点相减求和，交叉乘减求积。}
先只考虑横开口双曲线
$\dfrac{x^2}{a^2}-\dfrac{y^2}{b^2}=1$ 与非竖直直线 $y=kx+m$ 的\textbf{两个不同交点}。

\textbf{求和为什么用相减？} 因为 $x_1^2-x_2^2=(x_1-x_2)(x_1+x_2)$：
两点方程相减时，平方差直接带出所求的 $x_1+x_2$，而右端的 $1-1$ 自动消失。
再用直线关系 $y_1-y_2=k(x_1-x_2)$、
$y_1+y_2=k(x_1+x_2)+2m$，无需完整联立展开。

\textbf{求积为什么要交叉乘减？} 令
$F(x)=\dfrac{x^2}{a^2}-\dfrac{(kx+m)^2}{b^2}-1$，
两个交点使 $F(x_1)=F(x_2)=0$。
把第一个等式乘 $x_2$、第二个等式乘 $x_1$，然后相减：
\[
 x_2F(x_1)-x_1F(x_2)=0.
\]
此时展开平方所产生的\textbf{一次交叉项}都是 $2kmx_1x_2$，恰好互相抵消；
留下的平方项含 $x_1x_2(x_1-x_2)$，从而直接导出乘积。

\textbf{记忆口令：} \fbox{直接减，出和；交叉乘减，出积。}
这里的交叉乘减是理解乘积关系的巧法，不一定比考场直接看二次项系数和常数项更省笔画。
\end{explanationbox}
\clearpage

% ===== 再补充：两点法的初学者逐步理解（此前内容和 TikZ 均未修改）=====
\begin{explanationbox}
\textbf{深入理解一：为什么“两点相减”专门适合求和？}

\textbf{先别看双曲线，先看普通数字。}
例如 $7^2-3^2=(7-3)(7+3)$。
两个平方一相减，不用把 $49-9$ 当成主要计算目标；
我们真正得到的是“\textbf{两数之差}\,$\times$\,\textbf{两数之和}”。
因此，如果题目需要 $x_1+x_2$，看到 $x_1^2$ 与 $x_2^2$，
自然就应想到：\textbf{把两个平方相减，再使用平方差公式}。

\textbf{把这个想法搬到双曲线。} 两个点分别满足
\[
\underbrace{\frac{x_1^2}{a^2}-\frac{y_1^2}{b^2}=1}_{\text{点 }P},
\qquad
\underbrace{\frac{x_2^2}{a^2}-\frac{y_2^2}{b^2}=1}_{\text{点 }Q}.
\]
左式减右式，\textbf{右端先消成 $1-1=0$}：
\[
\frac{x_1^2-x_2^2}{a^2}-\frac{y_1^2-y_2^2}{b^2}=0.
\]
两个平方差分别拆开：
\[
\boxed{\frac{(x_1-x_2)(x_1+x_2)}{a^2}
=\frac{(y_1-y_2)(y_1+y_2)}{b^2}}.
\]
注意这一步不是在求根；\textbf{所需的两个“和”，已经直接露出来了}。

\textbf{还卡着两个“差”，怎么办？} 刚好两点又在同一直线
$y=kx+m$ 上：
\[
 y_1=kx_1+m,\qquad y_2=kx_2+m.
\]
两式相减时，截距 $m-m$ 又自动消失：
\[
 \boxed{y_1-y_2=k(x_1-x_2)}.
\]
用它替换上面的 $y_1-y_2$，两个不同交点保证 $x_1-x_2\ne0$，
就能约掉共同的“差”，只剩两个“和”：
\[
 \boxed{\frac{x_1+x_2}{a^2}=\frac{k(y_1+y_2)}{b^2}}.
\]
\textbf{一眼记住：} \fbox{曲线相减，平方变“差$\times$和”；直线相减，消掉那个“差”。}
\end{explanationbox}
\clearpage

\begin{explanationbox}
\textbf{深入理解二：如何从两个坐标和走到 $x_1+x_2$？}

上页已经得到
\[
 \frac{x_1+x_2}{a^2}=\frac{k(y_1+y_2)}{b^2}.
\]
但是出现了 $y_1+y_2$，别慌，\textbf{直接把两点的直线方程相加}：
\[
\begin{aligned}
 y_1+y_2
 &= (kx_1+m)+(kx_2+m)\\
 &=k(x_1+x_2)+2m.
\end{aligned}
\]
注意：\textbf{求差时 $m$ 抵消；求和时 $m$ 变成 $2m$}。
于是
\[
\begin{aligned}
 \frac{x_1+x_2}{a^2}
 &=\frac{k\,[k(x_1+x_2)+2m]}{b^2}\\
 &=\frac{k^2(x_1+x_2)}{b^2}+\frac{2km}{b^2}.
\end{aligned}
\]
把含 $x_1+x_2$ 的项挪到左侧：
\[
 \boxed{\left(\frac1{a^2}-\frac{k^2}{b^2}\right)
 (x_1+x_2)=\frac{2km}{b^2}}.
\]
如 $b^2-a^2k^2\ne0$，便能解出
\[
 \boxed{x_1+x_2=\frac{2a^2km}{b^2-a^2k^2}}.
\]
\textbf{这三步背后的目的：}
先用曲线相减\emph{制造}坐标和，再用直线之差\emph{消去差}，
最后用直线之和\emph{把 $y$ 的和换成 $x$ 的和}。
整个过程没有展开消元后的二次式。

\textbf{额外的几何理解。} 设弦中点为 $S(s,t)$，
则 $x_1+x_2=2s$、$y_1+y_2=2t$，
上面的关键等式也就是 $s/a^2=kt/b^2$。
所以这种“求和法”实际上是在找\textbf{弦的中点}。
\end{explanationbox}
\clearpage

\begin{explanationbox}
\textbf{深入理解三：为什么“交叉乘减”恰好求积？}

先观察一个\textbf{最简单的二次式}，不带双曲线参数：
\[
 ux^2+vx+w=0.
\]
因为 $x_1,x_2$ 都是根，所以
\[
 ux_1^2+vx_1+w=0\quad(1),\qquad
 ux_2^2+vx_2+w=0\quad(2).
\]
\textbf{目标是去掉 $vx_1$ 和 $vx_2$ 这两个一次项。}
若只是 $(1)-(2)$，留下 $v(x_1-x_2)$，它\emph{并未消失}。
但如果把 $(1)$ 乘 $x_2$，把 $(2)$ 乘 $x_1$：
\[
\begin{array}{rcl}
 (1)\times x_2: & ux_1^2x_2+\boxed{vx_1x_2}+wx_2=0,\\[3pt]
 (2)\times x_1: & ux_2^2x_1+\boxed{vx_2x_1}+wx_1=0.
\end{array}
\]
看中间的框：$vx_1x_2=vx_2x_1$。\textbf{它们完全一样！}
因此第一行减第二行，一次项自动归零：
\[
 u\,x_1x_2(x_1-x_2)+w(x_2-x_1)=0.
\]
又因为 $x_2-x_1=-(x_1-x_2)$，所以
\[
 (x_1-x_2)(ux_1x_2-w)=0.
\]
当 $x_1\ne x_2$ 时，得到
\[
 \boxed{ux_1x_2=w}.
\]
\textbf{这不是硬凑出来的交叉相乘。} 我们从“要消掉一次项”出发，
倒过来想：\emph{用什么乘数能让 $x_1$ 和 $x_2$ 变成同一个乘积？}
答案当然是第一个乘 $x_2$、第二个乘 $x_1$。

\textbf{迁移到双曲线时只有一处变化：}
把二次式中的 $u$ 换成 $1/a^2-k^2/b^2$，
把常数 $w$ 换成 $-1-m^2/b^2$。
下一页的详细证明会解释，这两个量如何直接从平方中读出。
\end{explanationbox}
\clearpage
